% GENERATED FILE. Edit canonical lessons, then run npm run generate:books.
% canonical-source-hash: fnv1a64:a2ef66588300807d
% canonical-lessons: HI-C413-saray, HI-C413-dharmsala, HI-C413-lift, HI-C413-thaharna, HI-C413-bichana, HI-R413-review-406-413

\chapter{Inns and Rest Houses}
\label{ch:hi-inns-and-rest-houses}
\hlchaptermodality{413}

\begin{chapteropening}
\textbf{By the end of this chapter:} \emph{I can stay at an inn or a rest house, take the lift and get a fresh sheet on the bed.}

Complete the last lesson of chapter 413: I can stay at an inn or a rest house, take the lift and get a fresh sheet on the bed.
\end{chapteropening}

\section[\texorpdfstring{\hi{सराय} (sarāy)}{सराय (sarāy)}]{\hi{सराय} (sarāy) — an inn}

\label{lesson:HI-C413-saray}



\begin{quote}
Before the new one: say the Hindi for to unload, then the Hindi for to put on board.
\end{quote}

\subsection*{You'll want to know: \hi{सराय}}
\textbf{\hi{सराय}} — \emph{sarāy} — \textquotedblleft{}an inn\textquotedblright{}.

\begin{quote}
\hi{पहले} \hi{मुसाफ़िर} \hi{सराय} \hi{में} \hi{रुकते} \hi{थे।}
\end{quote}

\begin{quote}
\emph{pahle musāfir sarāy meṁ rukte the.}
\end{quote}

\begin{quote}
In the old days, travellers used to stop at an inn.
\end{quote}

Feminine.

\begin{cousinweb}
\textbf{\hi{सराय}} (\emph{sarāy}) is usually traced to Persian \emph{sarāy}, \textquotedblleft{}a large house, a palace\textquotedblright{}, which also named the walled inns built along the old trade roads; many towns still have a lane or a quarter called \hi{सराय}.
\end{cousinweb}

\begin{grammarlens}[title={What you've built}]
One more word for getting around.
\end{grammarlens}

\subsection*{Your turn}
\begin{itemize}
\raggedright
  \item \emph{Say it:} \emph{sarāy}
  \item \emph{Say it:} \emph{sarāy}, once more
  \item \emph{Say it:} \emph{pahle musāfir sarāy meṁ rukte the.} — In the old days, travellers used to stop at an inn.
  \item \emph{Recall:} say the Hindi for to unload, then the Hindi for to put on board, then say \emph{sarāy} again
\end{itemize}

\begin{tcolorbox}[breakable,colback=teal!4,colframe=teal!35!black,title={Before you move on}]
What does \hi{सराय} mean? (\textquotedblleft{}An inn\textquotedblright{}.) Say it once more.
\end{tcolorbox}
\section[\texorpdfstring{\hi{धर्मशाला} (dharmśālā)}{धर्मशाला (dharmśālā)}]{\hi{धर्मशाला} (dharmśālā) — a rest house (for pilgrims)}

\label{lesson:HI-C413-dharmsala}



\begin{quote}
Before the new one: say the Hindi for to put on board, then the Hindi for an inn.
\end{quote}

\subsection*{You'll want to know: \hi{धर्मशाला}}
\textbf{\hi{धर्मशाला}} — \emph{dharmśālā} — \textquotedblleft{}a rest house (for pilgrims)\textquotedblright{}.

\begin{quote}
\hi{मंदिर} \hi{के} \hi{पास} \hi{एक} \hi{सस्ती} \hi{धर्मशाला} \hi{है।}
\end{quote}

\begin{quote}
\emph{mandir ke pās ek sastī dharmśālā hai.}
\end{quote}

\begin{quote}
There is a cheap rest house near the temple.
\end{quote}

Feminine.

\begin{cousinweb}
\textbf{\hi{धर्मशाला}} (\emph{dharmśālā}) is a tatsama, Sanskrit \emph{dharmaśālā} (Monier-Williams), from \emph{dharma}, \textquotedblleft{}duty, charity\textquotedblright{}, and \emph{śālā}, \textquotedblleft{}a hall\textquotedblright{}: a house kept for pilgrims and travellers at little or no cost.
\end{cousinweb}

\begin{grammarlens}[title={What you've built}]
One more word for getting around.
\end{grammarlens}

\subsection*{Your turn}
\begin{itemize}
\raggedright
  \item \emph{Say it:} \emph{dharmśālā}
  \item \emph{Say it:} \emph{dharmśālā}, once more
  \item \emph{Say it:} \emph{mandir ke pās ek sastī dharmśālā hai.} — There is a cheap rest house near the temple.
  \item \emph{Recall:} say the Hindi for to put on board, then the Hindi for an inn, then say \emph{dharmśālā} again
\end{itemize}

\begin{tcolorbox}[breakable,colback=teal!4,colframe=teal!35!black,title={Before you move on}]
What does \hi{धर्मशाला} mean? (\textquotedblleft{}A rest house (for pilgrims)\textquotedblright{}.) Say it once more.
\end{tcolorbox}
\section[\texorpdfstring{\hi{लिफ़्ट} (lifṭ)}{लिफ़्ट (lifṭ)}]{\hi{लिफ़्ट} (lifṭ) — a lift}

\label{lesson:HI-C413-lift}



\begin{quote}
Before the new one: say the Hindi for an inn, then the Hindi for a rest house.
\end{quote}

\subsection*{You'll want to know: \hi{लिफ़्ट}}
\textbf{\hi{लिफ़्ट}} — \emph{lifṭ} — \textquotedblleft{}a lift\textquotedblright{}.

\begin{quote}
\hi{लिफ़्ट} \hi{बंद} \hi{है}, \hi{सीढ़ी} \hi{से} \hi{जाइए।}
\end{quote}

\begin{quote}
\emph{lifṭ band hai, sīṛhī se jāiye.}
\end{quote}

\begin{quote}
The lift is out of order; please take the stairs.
\end{quote}

Feminine.

\begin{cousinweb}
\textbf{\hi{लिफ़्ट}} (\emph{lifṭ}) is borrowed from English \emph{lift}; it is feminine in Hindi.
\end{cousinweb}

\begin{grammarlens}[title={What you've built}]
One more word for getting around.
\end{grammarlens}

\subsection*{Your turn}
\begin{itemize}
\raggedright
  \item \emph{Say it:} \emph{lifṭ}
  \item \emph{Say it:} \emph{lifṭ}, once more
  \item \emph{Say it:} \emph{lifṭ band hai, sīṛhī se jāiye.} — The lift is out of order; please take the stairs.
  \item \emph{Recall:} say the Hindi for an inn, then the Hindi for a rest house, then say \emph{lifṭ} again
\end{itemize}

\begin{tcolorbox}[breakable,colback=teal!4,colframe=teal!35!black,title={Before you move on}]
What does \hi{लिफ़्ट} mean? (\textquotedblleft{}A lift\textquotedblright{}.) Say it once more.
\end{tcolorbox}
\section[\texorpdfstring{\hi{ठहरना} (ṭhaharnā)}{ठहरना (ṭhaharnā)}]{\hi{ठहरना} (ṭhaharnā) — to stop over; to stay (somewhere)}

\label{lesson:HI-C413-thaharna}



\begin{quote}
Before the new one: say the Hindi for a rest house, then the Hindi for a lift.
\end{quote}

\subsection*{You'll want to know: \hi{ठहरना}}
\textbf{\hi{ठहरना}} — \emph{ṭhaharnā} — \textquotedblleft{}to stop over; to stay (somewhere)\textquotedblright{}.

\begin{quote}
\hi{हम} \hi{दो} \hi{रात} \hi{एक} \hi{होटल} \hi{में} \hi{ठहरे।}
\end{quote}

\begin{quote}
\emph{ham do rāt ek hoṭal meṁ ṭhahre.}
\end{quote}

\begin{quote}
We stayed two nights at a hotel.
\end{quote}

Past (he …): \hi{वह} \hi{धर्मशाला} \hi{में} \hi{ठहरा}, \emph{vah dharmśālā meṁ ṭhahrā}.

\begin{cousinweb}
\textbf{\hi{ठहरना}} (\emph{ṭhaharnā}) goes back to \emph{*stabhira}, \textquotedblleft{}fixed, firm\textquotedblright{}, a form CDIAL reconstructs (CDIAL 13680), which lists Hindi \emph{ṭhaharnā}, \textquotedblleft{}to stop\textquotedblright{}. To ṭhaharnā is to stand still in one place for a while.
\end{cousinweb}

\begin{grammarlens}[title={What you've built}]
One more word for getting around.
\end{grammarlens}

\subsection*{Your turn}
\begin{itemize}
\raggedright
  \item \emph{Say it:} \emph{ṭhaharnā}
  \item \emph{Say it:} \emph{ṭhaharnā}, once more
  \item \emph{Say it:} \emph{ham do rāt ek hoṭal meṁ ṭhahre.} — We stayed two nights at a hotel.
  \item \emph{Recall:} say the Hindi for a rest house, then the Hindi for a lift, then say \emph{ṭhaharnā} again
\end{itemize}

\begin{tcolorbox}[breakable,colback=teal!4,colframe=teal!35!black,title={Before you move on}]
What does \hi{ठहरना} mean? (\textquotedblleft{}To stop over; to stay (somewhere)\textquotedblright{}.) Say it once more.
\end{tcolorbox}
\section[\texorpdfstring{\hi{बिछाना} (bichānā)}{बिछाना (bichānā)}]{\hi{बिछाना} (bichānā) — to spread out (a sheet or a mat); to make up (a bed)}

\label{lesson:HI-C413-bichana}



\begin{quote}
Before the new one: say the Hindi for a lift, then the Hindi for to stop over.
\end{quote}

\subsection*{You'll want to know: \hi{बिछाना}}
\textbf{\hi{बिछाना}} — \emph{bichānā} — \textquotedblleft{}to spread out (a sheet or a mat); to make up (a bed)\textquotedblright{}.

\begin{quote}
\hi{नौकर} \hi{ने} \hi{बिस्तर} \hi{पर} \hi{नई} \hi{चादर} \hi{बिछाई।}
\end{quote}

\begin{quote}
\emph{naukar ne bistar par naī cādar bichāī.}
\end{quote}

\begin{quote}
The servant spread a fresh sheet on the bed.
\end{quote}

Past, with \hi{ने} (he or she …): \hi{उसने} \hi{फ़र्श} \hi{पर} \hi{दरी} \hi{बिछाई}, \emph{usne farś par darī bichāī}.

\begin{cousinweb}
\textbf{\hi{बिछाना}} (\emph{bichānā}) goes back to \emph{*vicchādayati}, \textquotedblleft{}spreads out\textquotedblright{}, a form CDIAL reconstructs (CDIAL 11692), which lists Hindi \emph{bichānā}, \textquotedblleft{}to spread\textquotedblright{}; \hi{बिस्तर} \hi{बिछाना} is to make up a bed on the floor or a cot.
\end{cousinweb}

\begin{grammarlens}[title={What you've built}]
One more word for getting around.
\end{grammarlens}

\subsection*{Your turn}
\begin{itemize}
\raggedright
  \item \emph{Say it:} \emph{bichānā}
  \item \emph{Say it:} \emph{bichānā}, once more
  \item \emph{Say it:} \emph{naukar ne bistar par naī cādar bichāī.} — The servant spread a fresh sheet on the bed.
  \item \emph{Recall:} say the Hindi for a lift, then the Hindi for to stop over, then say \emph{bichānā} again
\end{itemize}

\begin{tcolorbox}[breakable,colback=teal!4,colframe=teal!35!black,title={Before you move on}]
What does \hi{बिछाना} mean? (\textquotedblleft{}To spread out (a sheet or a mat); to make up (a bed)\textquotedblright{}.) Say it once more.
\end{tcolorbox}
\section[(dialogue)]{Looking back: from Taxis and Tyres to Inns and Rest Houses}

\label{lesson:HI-R413-review-406-413}



\begin{quote}
Say the Hindi for a taxi and a tyre.
\end{quote}

\subsection*{Your turn}
Say these aloud:

\begin{itemize}
\raggedright
  \item a taxi, a tyre, a puncture, a helmet, a licence
  \item a highway, a toll, a checkpoint, a convoy, to wash away
  \item a mile, a kos, a distance, a path, to chase off
  \item an airport, an aircraft, a pilot, a flight, arrival
  \item a passport, a visa, security, customs, a fee
  \item weight, a delay (formal), a belt, cancelled, departure
  \item abroad, a land far from home, a frontier, to unload, to put on board
  \item an inn, a rest house, a lift, to stop over, to spread out
\end{itemize}

\begin{tcolorbox}[breakable,colback=teal!4,colframe=teal!35!black,title={Before you move on}]
A taxi? (\textbf{\hi{टैक्सी}}, \emph{ṭaiksī}.) A tyre? (\textbf{\hi{टायर}}, \emph{ṭāyar}.) A puncture? (\textbf{\hi{पंचर}}, \emph{pañcar}.) A helmet? (\textbf{\hi{हेलमेट}}, \emph{helmeṭ}.) A licence? (\textbf{\hi{लाइसेंस}}, \emph{lāisens}.) A highway? (\textbf{\hi{राजमार्ग}}, \emph{rājmārg}.) A toll? (\textbf{\hi{टोल}}, \emph{ṭol}.) A checkpoint? (\textbf{\hi{नाका}}, \emph{nākā}.) A convoy? (\textbf{\hi{क़ाफ़िला}}, \emph{qāfilā}.) To wash away? (\textbf{\hi{बहाना}}, \emph{bahānā}.) A mile? (\textbf{\hi{मील}}, \emph{mīl}.) A kos? (\textbf{\hi{कोस}}, \emph{kos}.) A distance? (\textbf{\hi{फ़ासला}}, \emph{fāslā}.) A path? (\textbf{\hi{राह}}, \emph{rāh}.) To chase off? (\textbf{\hi{भगाना}}, \emph{bhagānā}.) An airport? (\textbf{\hi{एयरपोर्ट}}, \emph{eyarporṭ}.) An aircraft? (\textbf{\hi{विमान}}, \emph{vimān}.) A pilot? (\textbf{\hi{पायलट}}, \emph{pāyalaṭ}.) A flight? (\textbf{\hi{उड़ान}}, \emph{uṛān}.) Arrival? (\textbf{\hi{आगमन}}, \emph{āgaman}.) A passport? (\textbf{\hi{पासपोर्ट}}, \emph{pāsporṭ}.) A visa? (\textbf{\hi{वीज़ा}}, \emph{vīzā}.) Security? (\textbf{\hi{सुरक्षा}}, \emph{surakṣā}.) Customs? (\textbf{\hi{कस्टम}}, \emph{kasṭam}.) A fee? (\textbf{\hi{शुल्क}}, \emph{śulk}.) Weight? (\textbf{\hi{वज़न}}, \emph{vazan}.) A delay (formal)? (\textbf{\hi{विलंब}}, \emph{vilamb}.) A belt? (\textbf{\hi{पेटी}}, \emph{peṭī}.) Cancelled? (\textbf{\hi{रद्द}}, \emph{radd}.) Departure? (\textbf{\hi{प्रस्थान}}, \emph{prasthān}.) Abroad? (\textbf{\hi{विदेश}}, \emph{videś}.) A land far from home? (\textbf{\hi{परदेस}}, \emph{pardes}.) A frontier? (\textbf{\hi{सरहद}}, \emph{sarhad}.) To unload? (\textbf{\hi{उतारना}}, \emph{utārnā}.) To put on board? (\textbf{\hi{चढ़ाना}}, \emph{caṛhānā}.) An inn? (\textbf{\hi{सराय}}, \emph{sarāy}.) A rest house? (\textbf{\hi{धर्मशाला}}, \emph{dharmśālā}.) A lift? (\textbf{\hi{लिफ़्ट}}, \emph{lifṭ}.) To stop over? (\textbf{\hi{ठहरना}}, \emph{ṭhaharnā}.) To spread out? (\textbf{\hi{बिछाना}}, \emph{bichānā}.)
\end{tcolorbox}
